\chapter{Theoretische Grundlagen}

\textit{Point-and-Click}-Spiele zeichnen sich dadurch aus, dass die Interaktion
mit der Spielwelt ausschließlich durch das Anklicken von Spielobjekten auf dem
Bildschirm stattfindet. Es wird also weitestgehend auf das Verwenden von
Tastatureingaben verzichtet.

Der Kunststil \textit{Pixel Art} ist der Grafik früher Computerspiele der
80er- und 90er-Jahre nachempfunden. Charakteristisch für ihn ist eine niedrige
Auflösung, wodurch einzelne Pixel deutlich voneinander unterscheidbar werden.
Weiterhin weisen Werke dieses Stils in der Regel auch ein geringes Spektrum
verwendeter Farben auf.

Das Spiel wurde in der Programmiersprache Java, Version 21 umgesetzt,
da diese Inhalt des Informatikunterrichts ist und somit von allen Gruppenmitglieder mindestens
ansatzweise beherrscht wird.

Ein essentielles Grundkonzept bei der Implementierung war die \textit{objektorientierte Programmierung},
also das Aufteilen des Programms in \textit{Klassen}. So können wiederverwendbare und vielseitig
einsetzbare Datenstrukturen entworfen sowie eine übersichtliche Struktur des Codes geschaffen werden.
Da Java durch und durch objektorientiert aufgebaut ist, findet sich somit ein weiterer Grund
für die Verwendung der Programmiersprache.

Eine Klasse ist grundsätzlich eine Sammlung von Variablen, Konstanten und Funktionen.
Häufig dient sie auch als Definition einer Datenstruktur, sodass \textit{Instanzen}
bzw. \textit{Objekte} von ihr erstellt werden können. Diese speichern dann
jeweils separate Werte für die in der Klasse definierten Variablen, auch \textit{Attribute}
oder \textit{Felder} genannt.

Ein Weiteres Konzept, welches im Spiel zum Einsatz kam, ist die \textit{Vererbung}.
Diese beschreibt das sogenannte \textit{Erweitern} einer Klasse: Eine separate
Klasse, die \textit{Kindklasse}, erweitert die \textit{Elternklasse}. Erstere
erhält somit Zugriff auf die Attribute der letzteren, kann diese aber auch
überschreiben oder eigene hinzufügen. Letztendlich definiert die Elternklasse
so eine Gemeinsame Schnittstelle an Variablen, Konstanten und Funktionen,
welche auch über die Instanzen der Kindklasse verwendet werden können. Dabei
ist entscheidend, dass vernachlässigt werden kann, ob es sich bei einer Instanz
um eine solche der Elternklasse oder einer Kindklasse handelt - verschiedene
Implementierungen einer Schnittstelle können einheitlich verwendet werden.

Ursprünglich waren das \textit{Point-and-Click}-Interaktionsmodell sowie
\textit{Pixel Art} das Ergebnis von Hardwarelimitierungen früher Computer,
welche noch nicht fähig waren, mit komplexeren Grafiken und Spielabläufeen
umzugehen. Im Fall des \textit{CZGame} wurden diese Gestaltungsentscheidungen
sowohl bewusst als auch aufgrund der begrenzten Arbeitszeit und Kenntnisstände
bzw. Kompetenzen der Gruppenmitglieder getroffen.
